\section{用圆弧近似代替任意曲线弧}
\markboth{85. 用圆弧近似代替任意曲线弧}{85. 用圆弧近似代替任意曲线弧}

要了解一小段曲线弧的弯曲程度，最好的方法是把它近似看成一小段圆弧，因为我们对于圆的弯曲情况早已了如指掌。

\noindent
\begin{minipage}[t]{0.44\textwidth}
  \vspace{0pt}
  \centering
  \begin{tikzpicture}[scale=0.88, >=Stealth, line cap=round, line join=round]
    % 坐标系
    \draw[->] (-0.45, 0) -- (2.65, 0) node[below, font=\footnotesize] {$x$};
    \draw[->] (0, -0.35) -- (0, 2.95) node[left, font=\footnotesize] {$y$};
    \node[below left, font=\footnotesize] at (0, 0) {$O$};
    % 曲率圆 Q(0.85, 1.48), r = 1.15
    \coordinate (Q) at (0.85, 1.48);
    \draw[thick] (Q) circle (1.15);
    \fill (Q) circle (1.2pt) node[left, font=\footnotesize] {$Q$};
    % 圆上两点 M (angle -52 deg), M1 (angle -14 deg)
    \coordinate (M) at ({0.85 + 1.15*cos(-52)}, {1.48 + 1.15*sin(-52)});
    \coordinate (M1) at ({0.85 + 1.15*cos(-14)}, {1.48 + 1.15*sin(-14)});
    \fill (M) circle (1.3pt) node[below right=0.5pt, font=\footnotesize] {$M$};
    \fill (M1) circle (1.3pt) node[right=2pt, font=\footnotesize] {$M_1$};
    % 半径 QM, QM1
    \draw[thin] (Q) -- node[pos=0.48, left=1pt, font=\footnotesize] {$r$} (M);
    \draw[thin] (Q) -- node[pos=0.52, above=1pt, font=\footnotesize] {$r$} (M1);
    % 曲线 C (经过 M 与 M1, 在 M 处严格沿 38 度相切, 在 M1 处严格沿 76 度相切)
    \draw[very thick, curveblue]
      (-0.15, 0.22)
      .. controls (0.55, 0.16) and ($(M) - ({0.42*cos(38)}, {0.42*sin(38)})$) .. (M)
      .. controls ($(M) + ({0.28*cos(38)}, {0.28*sin(38)})$) and ($(M1) - ({0.28*cos(76)}, {0.28*sin(76)})$) .. (M1)
      .. controls ($(M1) + ({0.36*cos(76)}, {0.36*sin(76)})$) and (1.68, 2.08) .. (1.22, 2.34);
    \node[font=\footnotesize] at (1.10, 2.38) {$C$};
    % M 处切线 MT: 过 M(1.558, 0.574)，方向 38 deg，交 x 轴于 A0(0.823, 0)
    % M1 处切线 M1T1: 过 M1(1.966, 1.202)，方向 76 deg，交 x 轴于 A1(1.666, 0)
    \coordinate (A0) at (0.823, 0);
    \coordinate (A1) at (1.666, 0);
    \coordinate (T) at ({1.558 + 0.92*cos(38)}, {0.574 + 0.92*sin(38)});
    \coordinate (T1) at ({1.966 + 0.52*cos(76)}, {1.202 + 0.52*sin(76)});
    \draw[thin, name path=tanM] (A0) -- (T) node[right, font=\footnotesize] {$T$};
    \draw[thin, name path=tanM1] (A1) -- (T1) node[above right=-1pt, font=\footnotesize] {$T_1$};
    \path[name intersections={of=tanM and tanM1, by=P}];
    \fill (P) circle (1.1pt) node[right=2pt, font=\footnotesize] {$P$};
    % 夹角 alpha 与 alpha1
    \coordinate (Xright) at (2.6, 0);
    \pic[draw, angle radius=9pt, angle eccentricity=1.55, font=\scriptsize, "$\alpha$"] {angle=Xright--A0--T};
    \pic[draw, angle radius=8pt, angle eccentricity=1.65, font=\scriptsize, "$\alpha_1$"] {angle=Xright--A1--T1};
  \end{tikzpicture}
  \captionof*{figure}{图\ 85.1}
\end{minipage}
\hfill
\begin{minipage}[t]{0.53\textwidth}
  \vspace{0pt}
  \hspace{2em}如图 85.1，把曲线 $C$ 在它的点 $M$ 附近的一小段曲线弧近似地看成是圆弧，其圆心是 $Q$，半径是 $r$。$M_1$ 是曲线 $C$ 在 $M$ 附近的点，它也在圆 $Q$ 上，那么
  \[
    QM = QM_1 = r.
  \]
  设曲线 $C$ 在点 $M$ 的切线为 $MT$，$MT$ 与 $x$ 轴夹角为 $\alpha$；在 $M_1$ 的切线为 $M_1 T_1$，$M_1 T_1$ 与 $x$ 轴夹角为 $\alpha_1$，记
\end{minipage}

\[
  \alpha_1 - \alpha = \Delta\alpha,
\]
又设切线 $MT$ 和 $M_1 T_1$ 的交点为 $P$，则
\[
  \angle MQM_1 = \angle TPT_1 = \alpha_1 - \alpha = \Delta\alpha.
\]
圆心角 $MQM_1$ 所对的弧 $\arc{MM_1}$，也是曲线 $C$ 的弧，其长度为 $\Delta s$，所以
\[
  r = \frac{\Delta s}{\Delta\alpha}.
\]
又设两点 $M_1$ 和 $M$ 的横坐标之差为 $\Delta x$，把上式的分子和分母同除以 $\Delta x$，得到
\[
  r = \frac{\dfrac{\Delta s}{\Delta x}}{\dfrac{\Delta\alpha}{\Delta x}}.
\]
上面的考虑只是近似的，存在一定误差。当所取的曲线弧越小，即 $\Delta x$ 越小时，误差也越小。令 $\Delta x \to 0$，取极限，就得到
\[
  r = \frac{\dfrac{ds}{dx}}{\dfrac{d\alpha}{dx}} = \frac{s'}{\alpha'},
\]
这里用撇号表示对 $x$ 的导数。因为 $\tg\alpha = y'$，所以 $\alpha = \arctg y'$，
\[
  \alpha' = (\arctg y')' = \frac{(y')'}{1 + (y')^2} = \frac{y''}{1 + y'^2}.
\]
另一方面，根据第 12 节，又有
\[
  s' = \sqrt{1 + y'^2},
\]
将 $s'$ 和 $\alpha'$ 的表达式相除，就得到
\[
  r = \frac{(1 + y'^2)^{3/2}}{y''}.
\]

根据曲线方程 $y = f(x)$，可以求出在曲线任一点 $M(x, y)$ 处的 $y'$ 和 $y''$ 的值，代入上式可以确定对应的圆半径 $r$。圆心 $Q$ 位于曲线在 $M$ 处的法线上，并且在曲线凹的一侧，所以圆心 $Q$ 的位置也能确定。这样一来，对于一条任意曲线 $C$，在它的任一点 $M$ 处就可确定一个圆 $Q$，用圆 $Q$ 在点 $M$ 附近的小段圆弧来近似代替曲线 $C$ 在 $M$ 附近的小段曲线弧。圆 $Q$ 叫做曲线 $C$ 在点 $M$ 的曲率圆，它的圆心叫做曲率中心，半径 $r$ 叫做曲率半径。$r$ 的倒数叫做曲率，用字母 $\varkappa$ 表示：
\[
  \varkappa = \frac{1}{r} = \frac{y''}{(1 + y'^2)^{3/2}}.
\]

术语“曲率”的意思，就是指弯曲的程度。我们从日常生活经验可以知道，把一根铅丝弯曲成圆弧，圆半径越小，弯曲得越厉害。由此可知，任意曲线 $C$ 在它的点 $M$ 处的曲率半径 $r$ 越小，即曲率 $\varkappa$ 越大，曲线在 $M$ 点附近就弯曲得越厉害。半径为 $r$ 的圆，曲率处处为 $\frac{1}{r}$。直线在每一点的曲率都是零。

算出了任意曲线在一点 $M$ 的曲率 $\varkappa$ 或曲率半径 $r$，就知道了曲线在 $M$ 点附近的弯曲程度；而如果知道曲率中心 $Q$ 的位置，就更能知道曲线在 $M$ 附近朝着什么方向弯曲。所以，由于曲率圆的帮助，我们对于任意曲线在它的各部分的弯曲情况，也能做到心中有数了。

有时曲线不是用直角坐标方程给出，而是用参数方程给出。如果参数恰好是切线与 $x$ 轴的夹角 $\alpha$，或者参数是弧长 $s$，那么利用下面的公式将会非常方便：
\[
  r = \frac{ds}{d\alpha}, \quad \varkappa = \frac{1}{r} = \frac{d\alpha}{ds}.
\]
公式的来源只需要看一看 153 页上面的式子就会明白。% 原书 153 页即本节前面的 $r = \frac{ds/dx}{d\alpha/dx} = \frac{s'}{\alpha'}$
